\documentclass{article}
\usepackage[utf8]{inputenc}
\usepackage[spanish]{babel}

\title{Memoria Nexus Gaming Club}
\author{Albert y Sergio}
\date{Curso 2026-2027}

\begin{document}

\maketitle

\section{Introducción, problema y objetivos}

\subsection{Nombre del proyecto y frase de venta}
Nuestro proyecto se llama Nexus Gaming Club. 
La frase de venta es: Infraestructura de red y sistemas de alto rendimiento para un edificio físico dedicado al gaming y eSports.

\subsection{Problema o necesidad}
Los clubs de videojuegos tienen varios problemas hoy en día:
\begin{itemize}
    \item No tienen red propia y dependen de plataformas externas.
    \item La red no está preparada para jugar con buen ping.
    \item No hay servidores propios de juego ni de voz.
    \item La información está repartida por varios sitios.
    \item No hay seguridad real (ni firewall, ni copias de seguridad).
\end{itemize}

\subsection{Usuarios}
Los usuarios que van a usar nuestro sistema son:
\begin{itemize}
    \item Administradores: Para gestionar usuarios y servidores.
    \item Moderadores: Para gestionar el contenido.
    \item Jugadores: Para acceder a los escritorios y torneos.
        \item Visitantes: Para ver la página web.
\end{itemize}

\subsection{Objetivos}
Los objetivos que queremos conseguir este curso son:
\begin{enumerate}
    \item Montar la red local y el wifi.
    \item Configurar los servidores (DHCP, DNS, etc.).
    \item Hacer una página web con HTML y CSS.
    \item Poner un firewall y hacer copias de seguridad.
    \item Documentar todo en latex y subirlo a GitHub.
\end{enumerate}

\section{Diseño de la Red}

Para la red del club, hemos diseñado una topologia en estrella. 
Utilizaremos un Router para la salida a Internet, un Switch principal, un Servidor para los servicios (DHCP, DNS, web y juego) y un Punto de Acceso Wifi para los portatiles y moviles.

El diseño lo hemos realizado en la herramienta Packet Tracer de Cisco. 
Para la red Wi-Fi hemos configurado el SSID "Nexus\_Wifi" con seguridad WPA2-PSK (AES) para que la conexión sea segura y nadie sin contraseña pueda conectarse.

\section{Configuración de Servicios Básicos}

Para la infraestructura de Nexus, hemos virtualizado un servidor con Windows Server 2025. Una vez instalado, le asignamos una dirección IP estatica (192.168.1.10) para que los servicios no pierdan la conexion.

\subsection{DHCP}
Hemos instalado el rol de DHCP y creado un ambito llamado "Rango-Nexus" que asigna direcciones IP desde la 192.168.1.100 hasta la 192.168.1.200 a los puestos de juego.

\subsection{DNS}
Hemos configurado el servidor DNS creando una zona de busqueda directa llamada "nexus.local" y un registro de host (A) para el servidor.

\subsection{Servidor de Archivos}
Para compartir recursos, hemos creado una carpeta en el servidor llamada "Nexus\_Documentos" y la hemos compartido en red con permisos de control total para el grupo de usuarios.

\section{Servicios de Red Avanzados}

\subsection{Servidor de Voz (Mumble)}
Para la comunicacion del club, hemos instalado el servidor de voz Mumble en Windows Server. Configuramos el servicio para que se inicie automaticamente, y probamos su funcionamiento conectando un equipo cliente desde fuera de la máquina virtual, con resultado exitoso.

\subsection{Servidor de Juego}
Para las partidas del club, hemos instalado un servidor de Minecraft. Instalamos Java, descargamos el archivo del servidor, aceptamos los términos de la licencia y lo ejecutamos.

\end{document}